\section{Who Escapes?}
\label{sec:escape}

If remoteness were destiny, Iceland, the Faroe Islands, Singapore and
Mauritius would be poor. They are not. Any account of the edge must explain
them, and \cref{tab:escape} lines them up against the median Pacific island
state.

\begin{table}[t]
\centering
\caption{\takeaway{The escapees are small but mostly not far, and all built
  institutions or hubs.} Selected indicators,
  \windowStart--\windowEnd{} averages unless stated; Pacific column is the
  median across Pacific small island states.}
\label{tab:escape}
\small\sffamily
% generated by scripts/export_stats.py
\begin{tabular}{@{}l*{5}{r}@{}}
\toprule
 & Iceland & Faroe Is. & Singapore & Mauritius & \makecell[r]{Pacific\\SIDS} \\
\midrule
Population (thousands) & 373 & 53 & 5{,}680 & 1{,}262 & 128 \\
Remoteness (km) & 5{,}854 & 5{,}750 & 9{,}844 & 10{,}568 & 11{,}565 \\
GDP per capita, PPP (\$) & 65{,}036 & 68{,}941 & 127{,}296 & 24{,}543 & 6{,}018 \\
Government effectiveness & 1.65 & --- & 2.08 & 0.92 & -0.06 \\
Shipping connectivity (2025) & 29 & 21 & 614 & 79 & 26 \\
International cables & 4 & 2 & 32 & 4 & 2 \\
Tertiary emigration, 2010 (\%) & 17 & --- & 10 & 63 & 28 \\
Tourism, \% of exports (2015--19) & --- & --- & 3 & 34 & 38 \\
\bottomrule
\end{tabular}

\source{World Bank WDI and WGI; CEPII; UNCTAD; TeleGeography; IAB. --- = not available.}
\end{table}

Three lessons stand out.

\paragraph{Lesson 1: most escapees are small, not far.} Iceland and the
Faroes are among the \emph{least} remote economies on Earth by GDP-weighted
distance: they sit between the two largest consumer markets in the North
Atlantic. Their populations are tiny (\popIcelandK{} and \popFaroesK{}
thousand), but \cref{sec:edge} found that size alone carries little income
penalty. Their success is a warning against treating all \enquote{peripheries}
alike---and partial evidence \emph{for} the paper's distinction between
distance and scale. The genuinely far edge, the South Pacific, has no
equivalent success story among its small states.

\paragraph{Lesson 2: turn position into rent.} Each escapee found something
the centre needed and could only get from them. Iceland and the Faroes
control rich fisheries and, in Iceland's case, cheap geothermal and hydro
power. Mauritius built export processing, financial services and tourism;
tourism alone earned \tourismMauritius\% of its exports before the pandemic.
Singapore made itself the hub that others must route through---the
opposite of the edge. Its shipping connectivity index is \lsciSingapore{} against a Pacific median
of \pacificMedianLSCI, and it has \cablesSingapore{} international cables.

\paragraph{Lesson 3: institutions do the heavy lifting.} All four score far
above the Pacific median on government effectiveness (Singapore \govSingapore,
Iceland \govIceland, Mauritius \govMauritius, against \govPacificMedian{}),
and in \cref{tab:regression} governance absorbed about half of the raw
remoteness penalty. Good institutions are partly a product of history,
including colonial history \parencite{acemoglu2001}, and cannot simply be
imported; but they are the lever over which small states have most control.

\begin{counterpoint}[Escape is not immunity]
Even escapees pay the talent tax. Mauritius had lost \emigMauritius\% of
its university graduates to the OECD by 2010, far more than Iceland
(\emigIceland\%) or Singapore (\emigSingapore\%), which can offer graduates
careers at home. Iceland's 2008 banking collapse showed how a small economy that
plugs into global finance can be overwhelmed by it. Escape is a trajectory,
not a permanent state.
\end{counterpoint}
